% Generated by scripts/make_paper_tables.py from the frozen results.
% Do not edit by hand; the next build overwrites it.
\begin{table}[tbp]
\centering
\caption{The note decomposed into the positions a desk would book, per \$1,000 of principal}
\label{tab:decomposition}
\small
\begin{tabular}{l r r r r r r}
\toprule
Note & Funding bond & Coupon strip & Worst-of put & Issuer's call & Without call & With call \\
\midrule
N1 & 778.77 & 299.87 & $-$135.72 & $-$27.99 & 942.92 & 914.93 \\
N2 & 778.77 & 291.16 & $-$151.68 & $-$22.34 & 918.25 & 895.91 \\
N3 & 864.65 & 234.00 & $-$94.29 & $-$40.09 & 1,004.37 & 964.27 \\
N4 & 913.12 & 171.01 & $-$67.22 & $-$36.38 & 1,016.92 & 980.54 \\
N5 & 920.50 & 90.38 & $-$102.08 & 51.99 & 908.81 & 960.80 \\
N6 & 913.12 & 194.24 & $-$119.14 & $-$34.23 & 988.22 & 954.00 \\
\bottomrule
\end{tabular}
\notetext{\textit{Notes:} The investor is long a funding bond and a coupon strip and short a worst-of put, and signs are from the investor's side. Where redemption is the issuer's decision the call is a cost to the investor and the column is negative. On the one auto-callable note it is positive: redemption is mechanical, returns principal early in states where the note would otherwise stay exposed, and so favours the holder. \textit{Without call} is the sum of the first three columns and \textit{With call} subtracts the fourth. The decomposition runs on 50,000 paths with no control variate, because every part has to be simulated under the same draws for the parts to add up; Table~\ref{tab:results} reports the headline estimator at 200,000 paths with the control variate. The last column is therefore a coarser estimate of the same quantity, not a competing one.}
\end{table}
